\documentclass[11pt,a4paper]{article}
\usepackage[T1]{fontenc}
\usepackage{lmodern,amsmath,amssymb,amsthm,mathtools,booktabs,microtype,xcolor}
\usepackage[margin=26mm]{geometry}
\usepackage[colorlinks=true,linkcolor=blue!45!black,citecolor=blue!45!black,urlcolor=blue!45!black]{hyperref}
\newtheorem{theorem}{Theorem}[section]
\newtheorem{lemma}[theorem]{Lemma}
\newtheorem{proposition}[theorem]{Proposition}
\theoremstyle{remark}\newtheorem{remark}[theorem]{Remark}
\newcommand{\Q}{\mathbb Q}\newcommand{\Z}{\mathbb Z}
\DeclareMathOperator{\Tr}{Tr}
\setlength{\emergencystretch}{2em}
\title{A family of quartics in 317 variables from aggregate trace heights}
\author{}\date{8 November 2026}
\begin{document}\maketitle
\begin{abstract}
We specify a family of integer quartics in 317 total variables by
291 quadratic equations and a complete coefficient table polynomial
in one fixed integer parameter. Under the arithmetic premises of
the specified source construction, every sufficiently large member
defines the integers in the rationals. The new ingredient is a
linear height comparison for the sum of four rational traces of
elliptic abscissas. It reduces three height-detector blocks to two
and replaces a field product by a rational sum. The absence of a parity restriction also removes all four odd-index certificates. An input relabeling retains exact integer extraction. The sufficient
parameter threshold is ineffective because the height comparison
has an ineffective additive constant. Thus this is an existence
theorem with an exact family specification; it does not determine
the coefficients of a particular defining member. The fully
specified 389-variable quartic remains a separate result.
\end{abstract}

\section{Statement and exact family}
Every variable count includes the input $X$. Put $\ell=80989$ and
\[
                         A_N=\ell^{2N+1}\qquad(N\ge1).
\]
The files \nolinkurl{quadratic-family317.json} and
\nolinkurl{polynomial-family317.json} specify 291 integer quadratic
polynomials $Q_{N,j}$ in 317 coordinates and
\[
                         P_N=\sum_{j=1}^{291}Q_{N,j}^2.
\]
The table uses powers of $A_N$ as a compact notation for integer
coefficients. Neither $A_N$ nor $N$ is a variable of $P_N$.

\begin{theorem}\label{thm:main}
Assume the arithmetic input contract of the specified quartic source
\cite[Section 2]{Quartic392}, including its quantitative weighted
prime-selection premise. There is an integer $N_0\ge27464$ such that,
for every integer $N\ge N_0$,
\[
 x\in\Z\quad\Longleftrightarrow\quad
 \exists\boldsymbol y\in\Q^{316},\ P_N(x,\boldsymbol y)=0.
\]
For every $N\ge1$, the polynomial has total degree exactly four,
integer content one, and coefficient height
\[
                         H(P_N)<2^{1520}\,80989^{4N+2}.
\]
For $N\ge27464$ it has exactly 5032 nonzero coordinate monomials.
The proof does not determine a numerical value of $N_0$.
\end{theorem}

The height comparison used below is a proved consequence of
Von Buhren's theorem, not a new arithmetic conjecture. Its
ineffectiveness prevents us from selecting an explicit defining
member from the family. Finite coefficient verification establishes
the specified polynomial identities; it cannot resolve this
effectiveness issue or prove the retained arithmetic premises.

\section{The aggregate trace}
We use the same curve and point as the source:
\[
 K=\Q(\rho),\quad \rho^2=2,\quad d=3-\rho,\qquad
 E:y^2=x^3-d^2x,
\]
\[
 P=2(1+\rho,2\rho)
   =\left(\frac{57-28\rho}{8},\frac{-660+417\rho}{32}\right).
\]
The retained root charts give four points
\[
 Q_1=[n]P,\quad Q_2=[a]P,\quad Q_3=[t]P,\quad Q_4=[b]P
\]
with nonzero integer labels, as proved below. Define the rational number
\[
 F=\frac12\Tr_{K/\Q}\bigl(x(Q_1)+x(Q_2)+x(Q_3)+x(Q_4)\bigr),
 \qquad B=\max(|n|,|a|,|t|,|b|).
\]
We write $h$ for absolute logarithmic Weil height and $\widehat h_E$
for canonical height on $E$.

\begin{lemma}\label{lem:square-chart}
Assume $2E(K)=\Z P$ and $P=2R_0$, where
$R_0=(1+\rho,2\rho)$. Every solution in $K$ of
\[
 x=h^2,\qquad j^2=x^2-d^2
\]
gives a point $Q=(x,hj)=[m]P$ for a nonzero integer $m$.
No parity condition on $m$ is imposed.
\end{lemma}
\begin{proof}
If $h=0$, then $j^2=-d^2$, impossible in the real field $K$.
Thus $x\ne0$, and the equations imply that $Q$ lies on $E$.
The standard squareclass homomorphism for this two-isogeny is
\[
 \delta_x:E(K)\longrightarrow K^\times/K^{\times2},\qquad
 \delta_x(Q)=[x(Q)]
\]
away from $O$ and $T_0=(0,0)$, with
$\delta_x(O)=1$ and $\delta_x(T_0)=[-d^2]$.
For completeness, its homomorphism identity follows from the
chord-and-tangent rule: the product of the three intersection
abscissas of a nonvertical line with the curve is the square of
its ordinate intercept. The special value at $T_0$ agrees with
$x(Q+T_0)=-d^2/x(Q)$; the vertical and tangent cases give the
same identity.

The group premise gives $2Q=mP=2mR_0$, and hence
$Q=mR_0+T$ for some $T\in E[2](K)$. The four torsion points
are $O,T_0,(d,0),(-d,0)$, with squareclasses
$1,-d^2,d,-d$. They all have uniform real signs at the two
embeddings. In contrast, $\delta_x(R_0)=[1+\rho]$ has mixed
real signs. Since $x(Q)$ is a nonzero square, $m$ must be even.
Then $\delta_x(T)$ must be a square. The negative values
$-d^2,-d$ are not squares, and $d$ is not a square because its
norm to $\Q$ is $7$. Therefore $T=O$ and $Q=[m/2]P$.
The index is nonzero because the point supplied by the chart is
finite. This also proves it is nontorsion.
\end{proof}

\begin{lemma}\label{lem:aggregate}
There are fixed constants $M_1>0$ and $M_2\ge1$ such that
\[
                         B^2\le M_1h(F)+M_2.
\]
Also $F>12$, so $1$ and $F^2$ are distinct positive rational squares.
No effective value of $M_2$ is claimed here.
\end{lemma}
\begin{proof}
The aggregate-trace height theorem \cite{Aggregate}
gives
\[
 \sum_{i=1}^4\widehat h_E(Q_i)\le M h(F)+M_0
\]
for fixed constants. The point $P$ has positive canonical height and
$\widehat h_E([j]P)=j^2\widehat h_E(P)$; rescaling proves the first
claim. The theorem applies to all nonzero integer labels and requires no parity condition. Its proof excludes every
positive-dimensional abelian translate through the indicated points
in a fiber of the aggregate map. Von Buhren's family height theorem
then supplies the linear inequality, with an ineffective additive
constant.

Each $Q_i$ is a nonzero double. Its abscissa is a positive square
at both real embeddings; the curve equation gives $x(Q_i)>d$ at
each embedding. Since $d$ and its conjugate sum to $6$, the rational
trace of each abscissa is greater than $3$. Summing gives $F>12$.
\end{proof}

\section{Two height detectors}
Use only the scales $1,F^2$. Every retained local detector has the
same point tests, formal logarithm approximations, and test orders
as the 389-variable construction \cite{Quartic389}. Its three
classes of test orders are $1,10,40$. In particular, its
nonvanishing is automatic at every consumed contact prime by the
local argument of that construction; no trace inverse witness is
reintroduced.

Let $H$ denote the retained contact-height constant and let
$M_0,M_1^*$ be the two contact radicals. We distinguish the radical
notation $M_1^*$ from the constant $M_1$ in
Lemma~\ref{lem:aggregate}. The following implication is the only
height consequence required of the detectors.

\begin{lemma}\label{lem:two-scales}
If $J\ne0$ is an integer and both $M_0^k$ and $(M_1^*)^k$
divide $J$, then
\[
 h(F)\le(H+\log3)|J|^{19/k},\qquad
 B^2\le C|J|^{19/k}+M_2,
 \quad C=M_1(H+\log3).
\]
\end{lemma}
\begin{proof}
The retained contact-height theorem bounds both $h(\lambda)$ and
$h(F^2\lambda)$ by $(H+\log3)|J|^{19/k}$. The factor $\log3$
accounts for the fractional linear change in the contact argument.
Since the detector domain excludes $\lambda=0$,
\[
 2h(F)=h(F^2)\le h(F^2\lambda)+h(\lambda).
\]
This proves the first inequality. Apply
Lemma~\ref{lem:aggregate} for the second.
\end{proof}

\section{Soundness for every sufficiently large parameter}
Put
\[
 \beta=51X^2+2504,\qquad c=|a-1|,
 \qquad
 \Delta=t^2-a(a-2)tb+(a-2)^2b^2-1.
\]
The scalar block uses
\[
 r=\beta(\beta-1)(\beta-2401)(32\beta+343)(2\beta-1),
 \qquad R=\frac{A_Nr}{1+\ell^2r},\qquad c_R=R+9792.
\]
All factors of $r$ are positive for rational $X$. The constant
$9792$ replaces the old shift $4896$; its prime factors are still
only $2,3,17$. The five excluded squareclasses at $\ell$ are
unchanged, because the old affine input $2X_{\rm old}+1$ is now
$X$. Thus the same Euler tests give
\[
 v_\ell(R)\in\{2N+1,2N-1\}.
\]
At every prime in $M_5(\beta)$, $R$ still has positive odd valuation
and $c_R$ is a unit. Consequently the local scalar argument,
including the exceptional-point exclusions, is unchanged.
Its integral group identity is
$[m](2G)=[4(a-1)]H$ and requires only $a-1\ne0$ integral,
not even. It gives, for every $N\ge1$,
\begin{equation}\label{eq:scalar}
 c\ge\ell^{L_N},\qquad
 L_N=\left\lfloor\frac{187(2N-1)}8\right\rfloor+1,
 \qquad h(\beta)\le Hc^{19/24}.
\end{equation}
These conclusions are proved before the final height comparisons;
they do not depend on the number of height detectors. In particular
$c$ tends uniformly to infinity with $N$, and $B\ge|a|\ge c-1$.
All exceptional elliptic-point exclusions and ratio-domain checks
from the scalar argument remain in force.

\begin{lemma}\label{lem:lucas-residual}
For every sufficiently large $N$, every rational solution of the
quadratic system satisfies $\Delta=0$.
\end{lemma}
\begin{proof}
If $\Delta\ne0$, the retained local comparisons make the fortieth
powers of both contact radicals divide $\Delta$. For $B\ge1$,
the displayed expression gives
\[
 |\Delta|\le2B^4+6B^3+5B^2+1\le14B^4.
\]
Lemma~\ref{lem:two-scales} would therefore imply
\[
 B^2\le C\,14^{19/40}B^{19/10}+M_2.
\]
Dividing by $B^2$ shows this is impossible once $B$ exceeds a
fixed threshold, since $19/10<2$. Equation~\eqref{eq:scalar}
ensures that threshold for every sufficiently large $N$.
\end{proof}

\begin{lemma}\label{lem:input}
For every sufficiently large $N$, every rational solution satisfies
$n=\beta$ and $X\in\Z$.
\end{lemma}
\begin{proof}
The elementary Lucas growth lemma in \cite[Section 5]{Quartic392}
states that $\Delta=0$, $b\ne0$, and $|a|>7$ imply
\[
                         \log B\ge\log|t|>|a|-3.
\]
Suppose $n-\beta\ne0$ and let $J$ be its nonzero reduced numerator.
The order-ten tests make both tenth powers of the contact radicals
divide $J$. Also $|J|\le2B\exp h(\beta)$. By
Lemma~\ref{lem:two-scales},
\[
 B^2\le C\,2^{19/10}B^{19/10}
                     \exp\!\left(\frac{19}{10}h(\beta)\right)+M_2.
\]
For $B^2\ge2M_2$ it follows that
\[
 \log B\le10\log(2C)+19\log2+19h(\beta)
 \le10\log(2C)+19\log2+19Hc^{19/24}.
\]
The last expression is $o(c)$ as $c\to\infty$, whereas
$|a|-3\ge c-4$. This contradicts the Lucas lower bound once
$c$ exceeds a fixed threshold. Equation~\eqref{eq:scalar}
provides that threshold for every sufficiently large $N$.

Thus $\beta=n$ is an integer. If $X=u/v$ in lowest terms
with $v>0$, then $v^2\mid51$. Since $51$ is squarefree, $v=1$,
so $X\in\Z$. No parity restriction on $n$ is used.
\end{proof}

The finitely many thresholds in these arguments depend on the
fixed constants $C,M_2,H$. Since $L_N$ increases with $N$, there
is one $N_0$ such that all conclusions hold for every $N\ge N_0$.
Increase it to be at least $27464$ if necessary. This proves the
soundness implication with the quantifiers asserted in
Theorem~\ref{thm:main}. The dependence on $M_2$ is precisely the
ineffective part.

\section{Completeness}
Fix $N\ge1$ and an integer input $X$. The unchanged scalar
normalizer and point construction admit the same prescribed finite
local depths and real approximation conditions as the source.
Increasing $N$ changes a fixed coefficient and finite depth
requirements; it does not add a rational variable or an
arithmetic existence premise. The source proof permits these
requirements to be fixed before the later contact primes are
chosen.

Here are the required small-prime conditions for the modified
label and shift. Directly from $\beta=51X^2+2504$ one obtains
\[
 v_2(r)=\begin{cases}3,&4\mid X,\\2,&4\nmid X,\end{cases}
 \qquad v_3(r)=1,\qquad r\equiv15\pmod{17}.
\]
For the first assertion, odd $X$ gives $\beta\equiv3\pmod8$,
with two valuation-one factors in $r$. If $X$ is even, the only
even factor is $\beta=4(51(X/2)^2+626)$, of valuation two or
three as displayed. At three the only divisible factor is
$2\beta-1=3(34X^2+1669)$, whose quotient is nonzero modulo
three. At seventeen, $\beta\equiv5$ and direct substitution
gives the stated residue.

The denominator $1+\ell^2r$ is a unit at two and three and is
$16$ modulo seventeen; $A_N$ is a unit at all three primes.
Since $9792=2^6\cdot3^2\cdot17$, this gives
\[
 v_2(9792/R)\ge3,\quad v_3(9792/R)=1,
 \quad v_{17}(9792/R)=1.
\]
Hence $c_R/R=1+9792/R$ is a square in each of
$\Q_2,\Q_3,\Q_{17}$. The scalar pole input $c_RRW^2$ is
therefore locally square, and the same simultaneous depth and
real approximation argument makes it lie in $(0,17)$ with a
square reduced denominator and all required local depths.

The input root index $n=\beta$ may now be even. The comparison
construction permits every nonzero integer label: its points
are $U_i=[hn_i]P$ and $V_i=[(h-1)n_i]P$, and the simultaneous
congruences for $h$ are independent of the $n_i$. The scalar
and Lucas choices for the other labels can retain their former
parities. Every such root has the square-abscissa chart of
Lemma~\ref{lem:square-chart}; the four additional odd-index
charts are no longer required.

After the four root points have been chosen, $F$ is a fixed
positive rational number. The two scales $1,F^2$ are distinct
positive rational squares by Lemma~\ref{lem:aggregate}.
The retained simultaneous prime-pattern and pole realization
argument applies to this fixed finite list. Coincident zero contact
offsets are the common pattern already present in the original
construction. The nonzero contact offsets from the two scales do
not collide: the ratio of any two nonzero elements of
\[
 \left\{-\tfrac12,-1,\tfrac1{2399},-\tfrac{32}{407}\right\}
\]
with the same sign has absolute value at most $407/32<144<F^2$.
All required avoidance congruences are incorporated in the same
finite-support selection. The quantitative separation argument
retains its stated weighted premise.

The resulting witnesses satisfy all retained comparisons, the two
detectors, and their normalized poles. Their quadratic gate values
are obtained by the same rational formulas; the local inverse
guards are omitted as in \cite{Quartic389}. The new scale coordinate
is $F^2/4$. This supplies a rational solution for every integer
input. Finally, a sum of rational squares vanishes exactly when
all its summands vanish. Together with the soundness argument,
this proves the projection assertion in Theorem~\ref{thm:main}.

\section{Exact coordinate accounting and verification}
Start with the specified 389-variable system, whose full source
hash is recorded in the generator. Delete the second component
height detector, including its pole certificate and its private
evaluation coordinates. This removes 52 coordinates and 50
quadratic rows. Of the nine field-product and scale coordinates,
retain only the coordinate formerly called
\nolinkurl{heightProduct.scale0.quadratic}; the other eight and
their defining rows are deleted. Rename the retained coordinate
$s_F$ and replace its one row by
\[
 (x_{A,0}^*+x_{S,0}^*+2x_{T,0}^*+x_{Z,0}^*)^2=4s_F.
\]
The starred symbols are literal stored coordinates. Their
combination on the left before squaring is exactly $F$; in
particular the factor two on the third coordinate is required
by its existing normalization. The retained first component
detector uses the actual input $4\lambda s_F=F^2\lambda$.
Every one of its downstream rows is unchanged. The full affine
detector normalization is independently verified in
\cite{Quartic389}.

Delete the three private odd-index coordinates and their three
rows for each of the four roots. There are no other occurrences
of these twelve coordinates. The retained square-abscissa charts
are sufficient by Lemma~\ref{lem:square-chart}.

Finally substitute into every retained row
\[
 X_{\rm old}=\frac{X-1}{2},\qquad
 Y_{\rm old}=\frac{Y-2X+1}{4},
\]
where $Y_{\rm old}$ was the stored square of the old input.
The old square row becomes $Y=X^2$, and
\[
 51(2X_{\rm old}+1)^2+2504=51X^2+2504.
\]
The transformation is affine and invertible on the full coordinate
space: $X=2X_{\rm old}+1$ and
$Y=4Y_{\rm old}+4X_{\rm old}+1$. Thus it preserves quadratic
degree. Replace the scalar shift $4896$ by $9792$ in its one
defining row. Clear each row's rational denominators and divide
its integer coefficient content; neither operation changes its
rational zero set.

The counts are therefore
\[
 389-52-8-12=317,\qquad361-50-8-12=291.
\]
The independent checker applies the inverse affine map to every
new row and recovers the corresponding old row, including the
changed scalar constant, up to a specified nonzero rational
multiple. It separately checks sole uses of every deleted
coordinate and the complete sum-of-squares coefficient table.
A parameter-independent coefficient one certifies content one
and degree exactly four. The small-prime congruences and the
five Euler tests at $80989$ are also recomputed independently.

There are 5047 terms when powers of $A_N$ are listed separately,
and 5032 coordinate monomials after collection. Every coefficient
is an integer polynomial in $A_N$ of degree at most two. The
largest sum of absolute coefficient-polynomial entries is less than
$2^{1520}$, giving the height bound in the theorem. For
$A_N>2^{1521}$, every nonconstant coefficient polynomial has its
leading term larger than the sum of the lower terms. Thus none
vanishes there; $N\ge27464$ more than suffices.

The complete reproduction sequence, in the construction directory, is
\begin{verbatim}
python3 build_quartic389.py
python3 check_quartic389.py
python3 verify_detector_normalizations.py
python3 build_quartic329.py
python3 check_quartic329.py
python3 build_quartic317.py
python3 check_quartic317.py
\end{verbatim}
The intermediate family is only a specified algebraic stage in
the construction. The exact final family and verification reports
distinguish the algebraic specification from the ineffective
existence threshold.

\begin{thebibliography}{9}
\bibitem{Quartic392} J.~Brox and contributors,
\emph{A forced prime condition and an explicit quartic in 392 variables},
8 November 2026, commit
\nolinkurl{de7aa60dd1794d7ea62a8a83cdc4bfd22380c74f},
\url{https://github.com/JoseBrox/diophantine-z-q_improvements}.
\bibitem{Quartic389} J.~Brox and contributors,
\emph{A quartic in 389 variables from local nonvanishing},
8 November 2026, accompanying proof and exact coefficient data.
\bibitem{Aggregate} J.~Brox and contributors,
\emph{Aggregate trace heights of elliptic multiples},
8 November 2026, accompanying proof.
\end{thebibliography}
\end{document}
